\chapter{离散随机变量}\label{ch:b2:randomvar}

随机变量以函数而非事件来组织概率计算。
在可数空间上，该理论依托
\cref{ch:b2:series}中的可和族：
期望是以样本空间为指标集的族之和，
其全部性质 --- 线性、变换公式、
独立变量的乘积公式 --- 都是关于可和族的定理。
本章证明马尔可夫、切比雪夫、柯西–施瓦茨和琴生的关键不等式，
并以经典分布与弱大数律作结，
其证明在有了切比雪夫不等式后
只需两行。

\section{随机变量及其分布}

\begin{definition}[离散随机变量；分布]\label{def:b2:randomvar:law}
设 $(\Omega, \P)$ 为可数概率空间。
称\emph{随机变量}\index{随机变量}为映射 $X \colon \Omega \to E$（$E$ 为任一
集合；当 $E = \R$ 时称为\emph{实}随机变量）。其\emph{分布}\index{分布!随机变量的}
（或\emph{分布}\index{分布}）是概率测度 $\P_X$，定义在
可数集 $X(\Omega)$ 上，由下式给出：
\[
\P_X(\{x\}) = \P(X = x)
= \P\bigl(\{\omega : X(\omega) = x\}\bigr) .
\]
\end{definition}

\begin{example}[经典分布]\label{ex:b2:randomvar:laws}
\leavevmode
\begin{itemize}
\item \emph{伯努利}分布 $\mathcal{B}(p)$：$X \in \{0, 1\}$，$\P(X =
1) = p$。某事件的示性变量。
\item \emph{二项}分布 $\mathcal{B}(n, p)$：$\P(X = k) =
\binom nk p^k(1-p)^{n-k}$，$0 \leq k \leq n$：$n$ 次独立伯努利试验中
成功的次数（高中分册；下文借助
独立变量之和重新证明）。
\item \emph{几何}分布 $\mathcal{G}(p)$：$\P(X = k) =
(1-p)^{k-1}p$，$k \in \N^*$：首次成功的
序号（见\cref{ex:b2:proba:geometric}）。
\item \emph{泊松}分布 $\mathcal{P}(\lambda)$：$\P(X = k) =
e^{-\lambda}\frac{\lambda^k}{k!}$，$k \in \N$ --- 由指数级数
可知其为概率测度。稀有事件的分布
（见\cref{ch:b2:genfun}）。
\end{itemize}
\end{example}

\begin{remark}[各分布分别刻画什么]
四种分布各回答一个基本问题：伯努利分布回答
“事情是否发生？”；二项分布回答“$n$ 次
试验中发生了几次？”；几何分布回答“第一次发生要等多久？”；
泊松分布回答“在给定速率下，试验很多而单次发生的概率很小时，
发生了多少次？”。认清所问的问题，
便是建模的九成功夫：指示变量之和对应二项分布，
等待时间对应几何分布，稀有事件的计数对应
泊松分布 --- 从二项分布到泊松分布的过渡，由稀有事件的分布律严格刻画，
见
\cref{ch:b2:genfun}。
\end{remark}

\begin{proposition}[几何
分布的无记忆性]\label{prop:b2:randomvar:memoryless}
若 $X \sim \mathcal{G}(p)$，则对所有 $m, n \in \N$：
\[
\P(X > m + n \mid X > m) = \P(X > n) ,
\]
且在 $\N^*$ 上，几何分布是唯一
具有此性质的分布。
\end{proposition}

\begin{proof}
对几何权重求和，得 $\P(X > n) = (1-p)^n$。于是
\[
\P(X > m + n \mid X > m)
= \frac{\P(X > m + n)}{\P(X > m)}
= \frac{(1-p)^{m+n}}{(1-p)^m} = (1-p)^n = \P(X > n).
\]
反之，若 $G(n) = \P(X > n)$ 满足 $G(m + n) =
G(m)G(n)$，$G(0) = 1$，则由归纳法得 $G(n) = G(1)^n$；$q =
G(1) \in \intco{0}{1}$，$q = 0$ 或分布为
$\mathcal{G}(1 - q)$：$\P(X = k) = G(k-1) - G(k) =
q^{k-1}(1 - q)$。
\end{proof}

\begin{example}[没有哪个结果“该出现了”]
掷一枚骰子，等待出现六点：等待时间 $X \sim
\mathcal G(1/6)$。无记忆性说明，在 $10$ 次
没有出现六点之后，\emph{剩余}等待时间 $X - 10$，
在 $X > 10$ 的条件下，仍服从 $\mathcal G(1/6)$，
即条件期望等待时间仍为 $6$ 次投掷，与开始时完全相同。
骰子不会“记得”，也没有哪次六点“该出现了” --- 赌徒
谬误是相信条件分布应已发生偏移。
反之，命题中关于唯一性的部分说明，
这种无差异性\emph{刻画}了几何
等待时间：任何预测从不更新的等待时间，
都服从几何分布。真实的排队与寿命通常会更新，
这正是人们借以发现它们并非
几何分布的途径。
\end{example}

\section{期望}

\begin{definition}[期望]\label{def:b2:randomvar:expectation}
称实随机变量 $X$ 在 $(\Omega, \P)$ 上
\emph{有期望}，当族
$\bigl(X(\omega)\,\P(\{\omega\})\bigr)_{\omega \in \Omega}$ 可和
（见\cref{ch:b2:series}）时，其\emph{期望}\index{期望}为
\[
\E(X) = \sum_{\omega \in \Omega} X(\omega)\,\P(\{\omega\}) .
\]
\end{definition}

\begin{theorem}[转移定理]\label{thm:b2:randomvar:transfer}
$X$ 有期望，当且仅当族 $\bigl(x\,\P(X =
x)\bigr)_{x \in X(\Omega)}$ 可和，此时
\[
\E(X) = \sum_{x \in X(\Omega)} x\,\P(X = x) .
\]
更一般地，设 $f \colon X(\Omega) \to \R$，则变量
$f(X)$ 有期望当且仅当 $\sum_x \abs{f(x)}\,\P(X = x) <
\infty$，此时 $\E(f(X)) = \sum_x f(x)\,\P(X = x)$。
\end{theorem}

\begin{proof}
将 $\Omega$ 划分为水平集 $\Omega_x = \{X = x\}$，$x
\in X(\Omega)$。由可和族的分组求和定理
（见\cref{ch:b2:series}），族
$(X(\omega)\P(\{\omega\}))_\omega$ 可和当且仅当每个分组都
（自动成立：$\sum_{\omega \in \Omega_x}\abs{x}\P(\{\omega\}) =
\abs x\,\P(X = x)$）\emph{且}分组和构成的族
$\bigl(x\,\P(X = x)\bigr)_x$ 可和 --- 此时总和
相等。对 $f(X)$，将已证明的结论用于变量 $Y =
f \circ X$，其水平集为 $\{Y = y\} =
\bigsqcup_{x : f(x) = y}\{X = x\}$；再作一次分组求和，
将 $\sum_y y\,\P(Y = y)$ 化为 $\sum_x
f(x)\,\P(X = x)$，此时分组将取值 $x$
按其像 $f(x)$ 归并，且一族绝对可和当且仅当
另一族绝对可和。
\end{proof}

\begin{theorem}[期望的性质]\label{thm:b2:randomvar:linearity}
在有期望的随机变量全体上：
\begin{enumerate}
\item （线性性）$\E(aX + bY) = a\,\E(X) + b\,\E(Y)$。
\item （非负性与单调性）$X \geq 0 \Rightarrow \E(X)
\geq 0$；$X \leq Y \Rightarrow \E(X) \leq \E(Y)$；以及
$\abs{\E(X)} \leq \E(\abs X)$。
\item （控制性）若 $\abs X \leq Z$ 且 $Z$ 有期望，
则 $X$ 也有期望。
\end{enumerate}
\end{theorem}

\begin{proof}
以上各条均是可和族之和的性质
（见\cref{ch:b2:series}）：和的线性、逐项保号，
以及可和性的控制判别法。（注意
线性性在 $\Omega$ 上的\emph{定义}中显然成立，
而在转移公式上则颇为麻烦 --- 这正是
预先定义 $\E$ 的好处之一）。
\end{proof}

\begin{example}\label{ex:b2:randomvar:expectations}
$X \sim \mathcal{B}(n, p)$：将 $X = X_1 + \dots + X_n$
写成伯努利指示变量之和，利用线性性得 $\E(X) = np$ ---
无需二项式系数。$X \sim \mathcal{G}(p)$：$\E(X) =
\sum_{k\geq1}k(1-p)^{k-1}p = p\cdot\frac{1}{(1 - (1-p))^2} =
\frac1p$，由收敛圆盘内的几何级数逐项求导得出
（见\cref{ch:b2:powerseries}）。$X \sim \mathcal{P}(\lambda)$：$\E(X)
= \sum_{k\geq1}k e^{-\lambda}\frac{\lambda^k}{k!} = \lambda
e^{-\lambda}\sum_{j\geq0}\frac{\lambda^j}{j!} = \lambda$。
\end{example}

\begin{example}[转移公式的应用]\label{ex:b2:randomvar:transferex}
设 $X \sim \mathcal P(\lambda)$，求
$\E\bigl(\frac1{1+X}\bigr)$ --- $\frac1{1+X}$ 的分布
本身较为繁琐，但转移公式从不需要它：
\[
\E\Bigl(\frac1{1+X}\Bigr)
= \sum_{k\geq0}\frac{1}{k+1}\,\eu^{-\lambda}
\frac{\lambda^k}{k!}
= \frac{\eu^{-\lambda}}{\lambda}\sum_{k\geq0}
\frac{\lambda^{k+1}}{(k+1)!}
= \frac{\eu^{-\lambda}}{\lambda}\bigl(\eu^\lambda - 1\bigr)
= \frac{1 - \eu^{-\lambda}}{\lambda} .
\]
两点启示。就计算而言：识别出平移后的指数级数
即是全部工作 --- 转移公式将
$f(X)$ 的期望化为级数运算。就结构而言：
朴素代入的值应为 $\frac1{1 + \E X} =
\frac1{1 + \lambda}$，而真实答案更大，
\[
\frac{1 - \eu^{-\lambda}}{\lambda} \geq
\frac{1}{1 + \lambda},
\]
正如琴生不等式对凸函数 $t
\mapsto \frac1{1+t}$ 的要求。凸函数作用后的期望
高于朴素代入的值，而转移公式
加上级数验证，使抽象的不等式具体化。
\end{example}

\begin{theorem}[独立性与乘积]\label{thm:b2:randomvar:product}
称随机变量 $X, Y$ 为\emph{独立}的，若 $\P(X = x, Y =
y) = \P(X = x)\P(Y = y)$ 对一切 $x, y$ 成立 --- 等价地，
事件 $\{X \in A\}$ 与 $\{Y \in B\}$ 对一切 $A,
B$ 均独立。若 $X$ 与 $Y$ 为独立的实随机变量，
且均有期望，则 $XY$ 有期望，并且
\[
\E(XY) = \E(X)\,\E(Y) .
\]
\end{theorem}

\begin{proof}
两种表述的等价性，由对
$(x, y) \in A \times B$ 上的逐点恒等式求和得出
（两次运用 $\sigma$ 可加性）。对乘积，二重族
$\bigl(xy\,\P(X = x)\P(Y = y)\bigr)_{(x,y)}$ 可和，因为
由族的富比尼定理（见\cref{ch:b2:series}）有
\[
\sum_{x, y}\abs x \abs y\,\P(X{=}x)\P(Y{=}y)
= \Bigl(\sum_x \abs x \P(X{=}x)\Bigr)
\Bigl(\sum_y \abs y \P(Y{=}y)\Bigr) < \infty ;
\]
由独立性，此族恰为 $\bigl(xy\,\P(X = x, Y =
y)\bigr)$，其和为 $\E(XY)$，此即转移公式用于
变量 $(X, Y) \mapsto xy$ 的结果；再由富比尼定理，
将这一非负和算得乘积 $\E(X)\E(Y)$。
\end{proof}

\begin{example}[乘积：有独立与无独立
两种情形]\label{ex:b2:randomvar:productcontrast}
掷两颗均匀骰子。若 $Y$ 为第二颗骰子（与第一颗
独立），则 $\E(XY) = \E(X)\E(Y) = 3.5^2 = 12.25$。若改为
$Y = X$（一颗骰子与自身的“乘积”），
\[
\E(X^2) = \frac{1 + 4 + 9 + 16 + 25 + 36}{6} = \frac{91}{6}
\approx 15.17 \neq 12.25 :
\]
两种情形的边缘分布相同，联合分布不同，
乘积的期望也不同。值得铭记的教训是：
$\E(XY)$ 是\emph{二元组}的泛函，而非两个
边缘分布的泛函 --- 而差 $\E(X^2) - \E(X)^2 \approx 2.92$
由柯尼希--惠更斯公式，恰为骰子的方差
$\frac{35}{12}$。
\end{example}

\section{方差、协方差与经典不等式}

\begin{definition}[矩与方差]\label{def:b2:randomvar:variance}
称 $X$ 具有\emph{二阶矩}，若 $X^2$ 有期望
（此时 $X$ 亦有期望，由控制：$\abs X \leq \frac{1 +
X^2}{2}$）。则其\emph{方差}\index{方差}与\emph{标准
差}\index{标准差}为
如下：
\[
V(X) = \E\bigl((X - \E(X))^2\bigr)
= \E(X^2) - \E(X)^2 ,
\qquad
\sigma(X) = \sqrt{V(X)} ,
\]
（第二式即\emph{柯尼希--惠更斯}公式 --- 由
展开平方并利用线性性：
\[
\E\bigl((X - \E X)^2\bigr)
= \E\bigl(X^2 - 2X\,\E X + \E(X)^2\bigr)
= \E(X^2) - 2\,\E(X)^2 + \E(X)^2 ,
\]
中间一项利用 $\E X$ 为常数）。对于具有二阶矩的 $X, Y$，
其\emph{协方差}\index{协方差}为
\[
\operatorname{Cov}(X, Y)
= \E\bigl((X - \E X)(Y - \E Y)\bigr)
= \E(XY) - \E(X)\E(Y) .
\]
\end{definition}

\begin{theorem}[方差的常用性质]\label{thm:b2:randomvar:variancerules}
对具有二阶矩的随机变量，有：
\begin{enumerate}
\item $V(aX + b) = a^2\,V(X)$；
\item $V(X + Y) = V(X) + V(Y) + 2\operatorname{Cov}(X, Y)$，
更一般地有
\[
V\Bigl(\sum_{i=1}^n X_i\Bigr)
= \sum_{i=1}^n V(X_i)
+ 2\sum_{i < j}\operatorname{Cov}(X_i, X_j) ;
\]
\item 若 $X, Y$ 独立，则 $\operatorname{Cov}(X, Y) = 0$
（逆命题不成立），故独立随机变量的方差
相加。
\end{enumerate}
\end{theorem}

\begin{proof}
\emph{1}与\emph{2}均由展开平方并利用线性性得到；
乘积 $X_iX_j$ 的期望由下文柯西--施瓦茨不等式
（或由 $\abs{X_iX_j} \leq \frac{X_i^2 + X_j^2}{2}$）得到。\emph{3}由
\cref{thm:b2:randomvar:product}应用于中心化后的
随机变量即得。逆命题的一个标准反例为：$X$ 在
$\{-1, 0, 1\}$ 上均匀分布，$Y = X^2$ 二者不相关
（$\E(XY) = \E(X^3) = 0 = \E X \cdot \E Y$），但显然不独立。
\end{proof}

\begin{theorem}[马尔可夫与切比雪夫
不等式]\label{thm:b2:randomvar:markov}
\leavevmode
\begin{enumerate}
\item （马尔可夫）若 $X \geq 0$ 有期望，则对每个 $a
> 0$：
\[
\P(X \geq a) \leq \frac{\E(X)}{a} .
\]
\item （切比雪夫）若 $X$ 具有二阶矩，则对每个
$\varepsilon > 0$：
\[
\P\bigl(\abs{X - \E(X)} \geq \varepsilon\bigr)
\leq \frac{V(X)}{\varepsilon^2} .
\]
\end{enumerate}
\end{theorem}

\begin{proof}
\emph{1.} 逐点地，$a\,\mathbf{1}_{X \geq a} \leq X$（在
事件上左端为 $a \leq X$；在其外则为 $0 \leq X$）。
两边取期望，得 $a\,\P(X \geq a) \leq \E(X)$，由单调性
及 $\E(\mathbf{1}_A) = \P(A)$。
\emph{2.} 对非负随机变量 $(X - \E
X)^2$ 取 $a = \varepsilon^2$ 应用马尔可夫不等式，事件 $\{(X - \E X)^2
\geq \varepsilon^2\}$ 恰为 $\{\abs{X - \E X} \geq
\varepsilon\}$。
\end{proof}

\begin{example}[不相关却
彼此紧扣]\label{ex:b2:randomvar:sumdiff}
掷两颗均匀骰子，$X$ 与 $Y$ 独立，并令 $S = X +
Y$，$D = X - Y$。由协方差的双线性，
\[
\operatorname{Cov}(S, D) = V(X) - V(Y) +
\operatorname{Cov}(Y, X) - \operatorname{Cov}(X, Y) = V(X) -
V(Y) = 0 :
\]
和与差不相关。它们独立吗？当然不
独立：$S = 12$ 迫使 $D = 0$，而 $\P(D = 0) =
\frac16$（无条件）。相关性只检验
\emph{线性}部分的依赖；这里的依赖由
约束 $S$ 与 $D$ 奇偶性相同所承载，
协方差无法察觉。（对此数对，零协方差
需要 $V(X) = V(Y)$：起作用的是同分布，
而非独立性。）
\end{example}

\begin{example}[马尔可夫不等式何时取等]\label{ex:b2:randomvar:markovsharp}
马尔可夫不等式取等号，当且仅当
估计 $a\,\mathbf 1_{X\geq a} \leq X$ 中没有任何浪费：
随机变量必须只取 $0$ 与 $a$ 两个值。具体地说，
若 $\P(X = a) = \pi$ 且 $\P(X = 0) = 1 - \pi$，则 $\E(X)
= a\pi$，且有
\[
\P(X \geq a) = \pi = \frac{\E(X)}{a} .
\]
一个现实的解读：在平均财富为
$100$、财富只能是 $0$ 或 $10^6$ 的人群中，
百万富翁的比例恰为 $10^{-4}$ --- 马尔可夫界恰好被取到，
这正是极端情形。只要 $X$ 取到中间值，
界便严格成立，往往差得很远；
但由极端情形可见，仅由均值
便不可能得到更好的不等式。
\end{example}

\begin{example}[切比雪夫不等式是精确的 --- 无需
进一步假设]\label{ex:b2:randomvar:chebsharp}
固定 $\varepsilon > 0$，$q \in \intoc01$，设 $X$ 取
$\pm\varepsilon$ 两值的概率各为 $\frac q2$，
取 $0$ 的概率为 $1 - q$。则 $\E(X) = 0$，$V(X) =
q\varepsilon^2$，且
\[
\P\bigl(\abs{X - \E X} \geq \varepsilon\bigr) = q
= \frac{V(X)}{\varepsilon^2} :
\]
切比雪夫不等式取等号。因此，仅用方差，不等式无法再改进
--- $1/\varepsilon^2$ 的衰减
正是二阶矩信息的精确代价。更快的衰减
需要更强的假设：随机变量的有界性
便得到\emph{指数型}集中性，正如
\cref{exo:b2:randomvar:7} 所预告，本章的周末问题
系统地展开了这一点。
\end{example}

\begin{theorem}[柯西--施瓦茨不等式与琴生不等式]\label{thm:b2:randomvar:jensen}
\leavevmode
\begin{enumerate}
\item （柯西--施瓦茨）若 $X, Y$ 有二阶矩，$XY$ 有
期望，且 $\E(XY)^2 \leq \E(X^2)\,\E(Y^2)$；由此
$\operatorname{Cov}(X,Y)^2 \leq V(X)V(Y)$。
\item （琴生）若 $\varphi \colon I \to \R$ 是
包含 $X(\Omega)$ 的区间上的凸函数，且 $X$、$\varphi(X)$ 有
期望，则
\[
\varphi\bigl(\E(X)\bigr) \leq \E\bigl(\varphi(X)\bigr) .
\]
\end{enumerate}
\end{theorem}

\begin{proof}
\emph{1.} $XY$ 的可和性：$\abs{XY} \leq \frac{X^2 + Y^2}2$。
映射 $(X, Y) \mapsto \E(XY)$ 是
二阶矩随机变量空间上的正对称双线性型，因此
\cref{ch:b2:quadratic} 中的抽象柯西--施瓦茨不等式适用
（对该不等式而言，正\emph{半}定即已足够）。
将其用于中心化随机变量，即得协方差的界。

\emph{2.} 首先，$m = \E(X)$ 属于 $I$：$I$ 是一个
包含 $X$ 全部取值的区间，期望具有单调性，故 $m$
介于 $\inf X(\Omega)$ 与 $\sup X(\Omega)$ 之间。由
凸函数的支撑线定理
（见\cref{ch:b2:realfun}），存在 $\alpha, \beta$ 使得
$\varphi(t) \geq \alpha t + \beta$ 对一切 $t \in I$ 成立，且
$\varphi(m) = \alpha m + \beta$。于是在 $\Omega$ 上逐点有
$\varphi(X) \geq \alpha X + \beta$；取期望得
\[
\E\bigl(\varphi(X)\bigr) \geq \alpha\,\E(X) + \beta
= \varphi\bigl(\E(X)\bigr). \qedhere
\]
\end{proof}

\begin{example}
取 $\varphi(t) = t^2$，琴生不等式给出 $\E(X)^2 \leq \E(X^2)$ ---
即方差非负；取 $\varphi(t) = 1/t$，其定义域为
$\intoo{0}{\infty}$ 上：$\frac{1}{\E X} \leq \E\bigl(\frac1X\bigr)$
--- 调和平均不超过算术平均，此处即
其随机形式。
\end{example}

\begin{remark}[常见陷阱]
(i) $\E(XY) = \E(X)\E(Y)$ \emph{要求}独立性（或至少
零协方差）。取 $Y = X$ 时，$\E(X^2) \neq
\E(X)^2$，只要 $V(X) > 0$。(ii) 同理，$V(X + X) =
4V(X)$，而非 $2V(X)$：方差只在独立
（或不相关）的加项之间才可相加。(iii) $\E(f(X))$ 不等于
$f(\E(X))$；对凸函数 $f$，琴生不等式甚至
可指出误差的方向，如
\cref{ex:b2:randomvar:transferex}。(iv) 存在性是一个真实的
假设：对圣彼得堡变量 $X = 2^K$，其中
$\P(K = k) = 2^{-k}$（$k \geq 1$），
\[
\sum_{k\geq1}2^k\cdot2^{-k} = \sum_{k\geq1}1 = \infty :
\]
$X$ 几乎必然有限，却没有期望，因而
该博弈不存在公平的入场价格。定义中 $\E$ 的可和性
并非形式上的吹毛求疵，而是
检测重尾的关键所在。(v) 最后，转移
转移定理要求\emph{绝对}可和，方可合法地
对取值上的求和作任何重排
（见\cref{ch:b2:series}）。
\end{remark}


\begin{example}[切比雪夫不等式用于一百次
抛掷]\label{ex:b2:randomvar:hundredtosses}
对 $X \sim \mathcal B(100, \frac12)$：$\E X = 50$，$V(X) =
25$。取 $\varepsilon = 6$，切比雪夫不等式给出
\[
\P(45 \leq X \leq 55) = \P(\abs{X - 50} < 6)
\geq 1 - \frac{25}{36} \approx 0.31 ,
\]
而精确的二项分布求和给出 $\approx 0.73$。
保证的 $31\%$ 远离真值，但它
\emph{仅}需要均值与方差 --- 同一结论
原样适用于任何满足 $\E = 50$、$V = 25$ 的
随机变量，无论多么奇特，且 \cref{ex:b2:randomvar:chebsharp} 表明
确有这样的变量使之取等。普适性是有代价的；
当分布确为二项分布时，
周末问题中的指数型工具可弥合大部分差距。
\end{example}

\begin{example}[一部分与其整体的
相关性]\label{ex:b2:randomvar:partwhole}
若 $X, Y$ 为独立同分布的随机变量，方差为
$\sigma^2 > 0$，则一个加项与和
$S = X + Y$ 的相关程度有多大？试计算
\[
\operatorname{Cov}(X, S) = \operatorname{Cov}(X, X) +
\operatorname{Cov}(X, Y) = \sigma^2 + 0 = \sigma^2,
\qquad V(S) = 2\sigma^2,
\]
于是相关系数为
\[
\rho(X, S) = \frac{\operatorname{Cov}(X,
S)}{\sigma(X)\,\sigma(S)}
= \frac{\sigma^2}{\sigma\cdot\sigma\sqrt2}
= \frac{1}{\sqrt2} \approx 0.707 ,
\]
无论共同分布为何——骰子、硬币与泊松计数皆然。若有
$n$ 个加项，同样的计算给出 $\rho(X_1, S_n) =
1/\sqrt n$：每个加项对总和的影响
按平方根规律衰减，这恰是涨落的
$\sqrt n$ 尺度在相关性上的投影。柯西–施瓦茨
不等式总保证 $\abs\rho \leq 1$；此处的界
恰在退化情形 $n = 1$ 时取等，
此后便按可预测的规律递减。
\end{example}

\begin{example}[由琴生不等式推出的加权
算术–几何平均不等式]\label{ex:b2:randomvar:amgm}
设 $Y$ 取正值 $a_1, \dots, a_k$，相应的
概率为 $\lambda_1, \dots, \lambda_k$。函数
$-\ln$ 在 $\intoo0\infty$ 上为凸函数，故由琴生不等式得
$-\ln\E(Y) \leq \E(-\ln Y)$，即
\[
a_1^{\lambda_1}a_2^{\lambda_2}\cdots a_k^{\lambda_k}
\;\leq\; \lambda_1a_1 + \lambda_2a_2 + \dots + \lambda_ka_k :
\]
即加权算术–几何平均不等式，其中等号
当且仅当 $Y$ 为常数时成立。等权重 $\lambda_i = \frac1k$
即得经典的算术–几何平均不等式。概率论悄悄证明了
一条纯代数定理：选取一个概率分布不过是
对凸组合的一种记账工具——
即\cref{ch:b2:affine}中的重心观点再次出现，此时以
琴生不等式为引擎。
\end{example}

\section{弱大数定律}

\begin{theorem}[弱大数定律]\label{thm:b2:randomvar:lln}
设 $(X_k)_{k \geq 1}$ 为两两独立的随机变量，
它们具有相同分布且存在二阶矩；记 $m = \E(X_1)$，
并记 $S_n = X_1 + \dots + X_n$。则对每个 $\varepsilon > 0$：
\[
\P\Bigl(\,\Bigl|\frac{S_n}{n} - m\Bigr| \geq \varepsilon\Bigr)
\;\leq\; \frac{V(X_1)}{n\,\varepsilon^2}
\xrightarrow[n \to \infty]{} 0 .
\]
\end{theorem}

\begin{proof}
由线性性得 $\E(S_n/n) = m$；由
\cref{thm:b2:randomvar:variancerules}（两两独立使协方差为零）得
$V(S_n) = n\,V(X_1)$，于是 $V(S_n/n) =
V(X_1)/n$。对 $S_n/n$ 应用切比雪夫不等式，即得
所需的界。
\end{proof}

\begin{remark}
本定理将概率与频率联系起来。取
$X_k$ 为独立重复试验中事件 $A$ 的示性变量，
则 $S_n/n$ 即为 $A$ 的观测频率，而大数定律
表明它在 $\P(A)$ 附近集中，其速率为
$\frac{p(1-p)}{n\varepsilon^2}$。\emph{强}大数定律（$S_n/n \to
m$，几乎必然）是本科三年级定理，其对四阶矩的证明
仍在可及范围之内：见\cref{exo:b2:randomvar:9}，
该题将博雷尔–坎特利引理用于切比雪夫型界。同样的切比雪夫估计
也支撑了魏尔斯特拉斯逼近定理的伯恩斯坦多项式证明，
该证明见\cref{ch:b2:funcseq}——其中的
计数引理\emph{正是}
伪装形式下的弱大数定律。
\end{remark}

\begin{example}[收集五十张赠券]
\label{ex:b2:randomvar:couponnumeric}
\cref{exo:b2:randomvar:3}中的赠券收集者，取 $n =
50$ 种互不相同的玩具，期望总数为
\[
\E(T_{50}) = 50\,H_{50} = 50\sum_{k=1}^{50}\frac1k
\approx 50 \times 4.499 \approx 225
\]
盒——是朴素猜测 $50$ 的四点五倍。
调和增长便是全部要点：前 $25$ 件玩具
大约在 $50\ln2 \approx 35$ 盒内集齐，而
\emph{最后一件}玩具单独平均需要 $50$ 盒（
参数为 $\frac1{50}$ 的几何等待）。完成型
问题由其收尾阶段主导，这正是
\cref{exo:b2:randomvar:12}得出阶为 $n$ 的涨落
——即最后那段几何等待的规模——围绕
均值 $n\ln n$ 附近。
\end{example}

\begin{example}[$n$ 至少要多大？]\label{ex:b2:randomvar:llnnumeric}
要使观测频率在 $\varepsilon = 0.01$ 之内逼近
$\P(A)$，置信度为 $95\%$，则切比雪夫界要求
\[
\frac{p(1-p)}{n\varepsilon^2} \leq \frac{1}{4n\varepsilon^2}
\leq 0.05,
\qquad\text{即}\qquad
n \geq \frac{1}{4\cdot0.05\cdot(0.01)^2} = 50\,000 .
\]
对 $\varepsilon$ 的依赖十分苛刻（二次），而
对置信度的依赖则较为温和（关于 $1/\alpha$ 线性）。两者皆
是\emph{界}本身的性质，而非真实情况的性质：周末问题中的
指数型不等式将置信度的代价由
$1/\alpha$ 降为 $\ln(1/\alpha)$——相同的规格
在那里约需 $18\,500$ 个样本——
而 $1/\varepsilon^2$ 的尺度是真实的，
且不可改进。知道界的哪一部分偏松，其价值
与界本身同样重要。
\end{example}

\begin{omfigure}
\begin{tikzpicture}[x=5.5cm, y=1.05cm]
  \draw[->] (-0.05,0) -- (1.1,0) node[below] {$S_n/n$};
  \draw[->] (0,-0.1) -- (0,2.35);
  \draw[dashed] (0.4,0) -- (0.4,2.15);
  \draw[dashed] (0.6,0) -- (0.6,2.15);
  \node[below, font=\scriptsize] at (0.4,0)
    {$m{-}\varepsilon$};
  \node[below, font=\scriptsize] at (0.5,0) {$m$};
  \node[below, font=\scriptsize] at (0.6,0)
    {$m{+}\varepsilon$};
  \draw[omDef, thick] plot[domain=0.06:0.94, samples=80]
    (\x, {0.8*exp(-(\x-0.5)^2/0.0288)});
  \draw[omThm, very thick] plot[domain=0.32:0.68, samples=80]
    (\x, {2.0*exp(-(\x-0.5)^2/0.0032)});
  \node[omDef, font=\small] at (0.16, 0.85) {$n$ 较小};
  \node[omThm, font=\small] at (0.71, 1.95) {$n$ 较大};
\end{tikzpicture}

{\small 大数定律的图示：
$S_n/n$（示意画出）的分布中心 $m$ 保持不变，但
随 $n$ 增大而变窄，于是落在
区间 $\intcc{m-\varepsilon}{m+\varepsilon}$ 之外的概率——即两条尾部——
趋于零。切比雪夫不等式将尾部
界为 $V(X_1)/(n\varepsilon^2)$；周末问题表明它们
实际上是指数级小的。}
\end{omfigure}

\begin{remark}[本册中的视角]
向前看，本章全部内容都服务于\cref{ch:b2:genfun}：
$\E(t^X)$ 是对 $X$ 的某个巧妙函数取的期望，
将整个分布打包成一个幂级数，矩则化为
 $1$ 处的导数，而随机和的瓦尔德型恒等式
支撑分支过程理论；独立变量的
乘积定理化为母函数的
乘法性。向后看，期望是以概率为权的
重心（见\cref{ch:b2:affine}），琴生
不等式是凸函数的支撑线几何
（见\cref{ch:b2:realfun}），而本章周末问题中的
指数矩方法即是对
 $\eu^{tX}$ 应用马尔可夫不等式——这是一个不等式，经一次恰当的
变量替换而推广，横跨三章。
\end{remark}

\section{练习}

\begin{exercise}[$\star$]\label{exo:b2:randomvar:1}
计算 $\E(X)$ 与 $V(X)$，其中 $X \sim \mathcal{B}(n, p)$（借助
指示变量），$X \sim \mathcal{P}(\lambda)$（证明 $V(X) =
\lambda$），$X \sim \mathcal{G}(p)$（证明 $V(X) =
\frac{1-p}{p^2}$；利用 $\E(X(X-1))$ 及
几何级数的二阶导数）。
\end{exercise}

\begin{exercise}[$\star$]\label{exo:b2:randomvar:2}
设 $X \sim \mathcal{P}(\lambda)$ 与 $Y \sim \mathcal{P}(\mu)$
相互独立。证明 $X + Y \sim \mathcal{P}(\lambda + \mu)$
（利用权的卷积与二项式定理），并证明
$X$ 的条件分布在 $X + Y = n$ 之下是二项分布
$\mathcal{B}\bigl(n, \frac{\lambda}{\lambda + \mu}\bigr)$。
\end{exercise}

\begin{exercise}[$\star$]\label{exo:b2:randomvar:3}
（优惠券收集者问题，期望）某麦片品牌在每盒中均匀地藏入
$n$ 件互不相同的玩具之一。设 $T_n$ 为
集齐全部 $n$ 件玩具所需的盒数。将 $T_n$ 写成
独立几何变量之和（即见到一件\emph{新}玩具的等待时间，当仍缺
$k$ 件时），证明
\[
\E(T_n) = n\sum_{k=1}^{n}\frac{1}{k} \sim n\ln n
\]
（由级数–积分比较得等价，
见\cref{ch:b2:comparison}）。
\end{exercise}

\begin{exercise}[$\star\star$]\label{exo:b2:randomvar:4}
设 $X \geq 0$ 为整值随机变量。证明\emph{尾公式}
\[
\E(X) = \sum_{n=1}^{\infty} \P(X \geq n)
\]
（当其中一侧有限时），通过写出 $X =
\sum_{n\geq1}\mathbf{1}_{X \geq n}$ 并交换求和次序
（对非负族应用富比尼定理）。由此对几何分布验证
$\E(X) = \frac1p$。
\end{exercise}

\begin{exercise}[$\star\star$]\label{exo:b2:randomvar:5}
（不放回抽样更为集中）一个罐子中有 $N$
个球，其中 $M$ 个为白色。不放回地抽取 $n \leq N$ 个球，
并令 $X$ 为白球的个数（\emph{超几何}分布）。利用
指示变量 $X = \sum_{i=1}^n Y_i$，其中 $Y_i$ 为第 $i$ 次抽取：
证明每个 $Y_i$ 服从参数为 $p = M/N$ 的伯努利分布（利用对称性！），
由此得 $\E(X) = np$（与有放回时完全相同），并证明
$\operatorname{Cov}(Y_i, Y_j) = -\frac{p(1-p)}{N-1} < 0$，其中 $i
\neq j$，从而 $V(X) = np(1-p)\frac{N - n}{N - 1} \leq np(1-p)$。
\end{exercise}

\begin{exercise}[$\star\star$]\label{exo:b2:randomvar:6}
设 $X$ 具有二阶矩。证明 $c \mapsto \E\bigl((X -
c)^2\bigr)$ 恰在 $c = \E(X)$ 处取最小值，最小值为
$V(X)$。再证明 $\P(X = \E(X)) = 1$ 当且仅当 $V(X) =
0$。\emph{（关于第二点：若 $V(X) = 0$，利用切比雪夫不等式，
取 $\varepsilon = 1/n$，并利用单调连续性，
见\cref{thm:b2:proba:continuity}。）}
\end{exercise}

\begin{exercise}[$\star\star\star$]\label{exo:b2:randomvar:7}
（集中性胜过马尔可夫不等式）设 $S_n \sim \mathcal{B}(n,
\frac12)$（$n$ 次公平投掷中正面出现的次数）。比较由
马尔可夫不等式给出的界（$\P(S_n \geq \frac{3n}{4})$）、切比雪夫不等式，
以及指数型（切尔诺夫）方法给出的界：
\[
\P\Bigl(S_n \geq \frac{3n}4\Bigr)
\leq \E\bigl(e^{tS_n}\bigr)e^{-3nt/4}
= \Bigl(\frac{1 + e^t}{2}\Bigr)^n e^{-3nt/4}
\quad (t > 0),
\]
并对 $t$ 求优化，得到指数级小的界。\emph{（当
$t = \ln 3$ 时，界为 $\bigl(2\cdot 3^{-3/4}\bigr)^n \approx
(0.877)^n$。）}
\end{exercise}

\begin{exercise}[$\star\star\star$]\label{exo:b2:randomvar:8}
（用概率方法再证魏尔斯特拉斯定理）设 $f \colon [0,1] \to \R$
为连续函数，且 $S_n \sim \mathcal{B}(n, x)$。证明
伯恩斯坦多项式 $B_nf(x) = \sum_{k=0}^n f\bigl(\frac
kn\bigr)\binom nk x^k(1-x)^{n-k}$ 等于
$\E\bigl[f\bigl(\frac{S_n}{n}\bigr)\bigr]$，并用概率语言重新推导
在\cref{ch:b2:funcseq}中的估计 $\abs{B_nf(x) - f(x)} \leq \omega_f(\delta) +
\frac{2\norm f_\infty}{4n\delta^2}$，
按 $\bigl|\frac{S_n}{n} -
x\bigr| \geq \delta$ 分情形，并用切比雪夫不等式。
\end{exercise}

\begin{exercise}[$\star\star\star$]\label{exo:b2:randomvar:9}
（四阶矩条件下的强大数定律）设 $(X_k)$ 为独立的，
同分布、均值为零（$\E X_1 = 0$），且
$\E(X_1^4) < \infty$。展开 $\E(S_n^4)$，并数出
留下的项（只有 $\E(X_i^4)$ 与 $\E(X_i^2X_j^2)$ 形式的项，$i
\neq j$），证明 $\E(S_n^4) \leq C n^2$，其中 $C$ 为常数。
由此推出 $\sum_n \P\bigl(\abs{S_n/n} \geq \varepsilon\bigr) < \infty$
对每个 $\varepsilon > 0$ 成立（用四阶马尔可夫不等式），并
由博雷尔--坎特利引理（见\cref{thm:b2:proba:borelcantelli}）得
$S_n/n \to 0$ 几乎必然成立，以适当的形式表述为：
事件 $\bigcap_{j}\bigcup_N\bigcap_{n \geq N}\{\abs{S_n/n} <
\frac1j\}$ 的概率为 $1$。
\end{exercise}

\begin{exercise}[$\star$]\label{exo:b2:randomvar:10}
掷两枚均匀骰子，令 $M$ 为两次掷出的点数中
较大者。利用\cref{exo:b2:randomvar:4}中的尾概率公式
（有限形式），证明
\[
\E(M) = \sum_{k=1}^{6}\P(M \geq k)
= 6 - \sum_{j=0}^5\Bigl(\frac j6\Bigr)^2 = \frac{161}{36}
\approx 4.47 .
\]
\end{exercise}

\begin{exercise}[$\star\star$]\label{exo:b2:randomvar:11}
设 $F_n$ 为 
$\{1, \dots, n\}$ 上均匀随机置换的不动点个数（$n \geq 2$）。记 $F_n =
\sum_i\mathbf 1_{\sigma(i) = i}$，计算 $\E(F_n) = 1$，
$\operatorname{Cov}(\mathbf 1_{\sigma(i)=i}, \mathbf
1_{\sigma(j)=j}) = \frac1{n^2(n-1)}$，当 $i \neq j$ 时，并
得 $V(F_n) = 1$：平均而言恰有一封信回到原位，
方差恰为 $1$，无论 $n$ 为何值。
\end{exercise}

\begin{exercise}[$\star\star\star$]\label{exo:b2:randomvar:12}
（优惠券收集问题，集中性）在
\cref{exo:b2:randomvar:3}的设定下，证明
\[
V(T_n) = \sum_{k=1}^n\frac{1 - k/n}{(k/n)^2}
\leq n^2\sum_{k=1}^n\frac{1}{k^2} \leq \frac{\pi^2}{6}n^2,
\]
利用几何分布各阶段的独立性以及 $V(\mathcal
G(p)) = \frac{1-p}{p^2}$（见\cref{exo:b2:randomvar:1}；
$\pi^2/6$ 的值见\cref{ex:b2:fourier:basel}）。
由切比雪夫不等式可得 $\dfrac{T_n}{n\ln n} \to 1$ \emph{依概率}：
收集者的总时间为 $n\ln n$，
精确到阶为 $n$ 的涨落。
\end{exercise}

\section{问题：集中性工具，从马尔可夫到
霍夫丁}

\begin{problem}[{周末问题 --- 指数型集中性的手算，
以及民意调查须询问多少人}]
\label{pb:b2:randomvar:1}
马尔可夫不等式只需一阶矩，换来 $1/a$ 的衰减；
切比雪夫不等式需要二阶矩，换来 $1/\varepsilon^2$ ---
而\cref{ex:b2:randomvar:chebsharp}表明，这就是这些矩所能换来的全部。
本题继续攀登余下的阶梯：
指数型（切尔诺夫）方法，及其针对掷硬币的\emph{精确}速率，
适用于一切有界变量的霍夫丁不等式，
收获 --- 明确而诚实的样本量，
用于民意调查、选举预测与硬币检验。
\index{集中性不等式}\index{霍夫丁不等式
}\index{切尔诺夫界}全文中，$S_n \sim
\mathcal B(n, p)$ 是 $n$ 个独立伯努利
变量之和，$\widehat p_n = S_n/n$ 为经验频率。

\medskip
\noindent\textbf{第 I 部分 --- 公平硬币上的校准。}
这里 $p = \frac12$，且 $a \in \intoo{\frac12}{1}$。
\begin{enumerate}
  \item 对水平 $an$ 用马尔可夫不等式，证明 $\P(S_n \geq an) \leq
        \frac1{2a}$，这一界甚至不趋于 $0$。
        马尔可夫不等式在哪里损失了这么多？
  \item 切比雪夫：利用公平二项分布的
        关于 $n/2$ 的对称性，证明
        \[
        \P(S_n \geq an) = \tfrac12\,
        \P\bigl(\abs{S_n - \tfrac n2} \geq n(a -
        \tfrac12)\bigr)
        \leq \frac{1}{8n(a - 1/2)^2},
        \]
        即 $\frac2n$（$a = \frac34$ 时），
        终究只是多项式衰减。
  \item （切尔诺夫，一般水平）计算
        $\E(\eu^{tS_n}) = \bigl(\frac{1 + \eu^t}2\bigr)^n$，并
        优化 $\P(S_n \geq an) \leq
        \E(\eu^{tS_n})\eu^{-tan}$ 关于 $t > 0$，证明
        最优的 $t$ 为 $\ln\frac{a}{1-a}$，并得
        \[
        \P(S_n \geq an) \leq \eu^{-n\,I(a)},
        \qquad
        I(a) = \ln 2 + a\ln a + (1-a)\ln(1-a) > 0 .
        \]
        验证当 $a = \frac34$ 时，所得界恢复为
        $\bigl(2\cdot3^{-3/4}\bigr)^n$，即
        \cref{exo:b2:randomvar:7}中的界。
  \item （指数是精确的）设 $k = an$ 为整数。
        由 $\binom nk a^k(1-a)^{n-k}$ 是
        某一概率分布的 $n + 1$ 项中的最大者这一事实，
        证明 $\binom nk \geq
        \frac{\eu^{nH(a)}}{n+1}$，其中 $H(a) = -a\ln a -
        (1-a)\ln(1-a)$，并由此推出匹配的下界
        \[
        \P(S_n \geq an) \geq \binom{n}{an}2^{-n}
        \geq \frac{\eu^{-n\,I(a)}}{n + 1} .
        \]
  \item 在 $n = 100$、$a =
        \frac34$ 时列表给出三个界：马尔可夫 $\frac23$，切比雪夫 $0.02$，
        切尔诺夫 $\approx 2.1\cdot10^{-6}$（真值为
        $\approx 2.8\cdot10^{-7}$）。试用一句话概括其中的启示。
\end{enumerate}

\medskip
\noindent\textbf{第二部分 --- 霍夫丁不等式。}
\begin{enumerate}[resume]
  \item （拉德马赫情形）设 $\varepsilon = \pm1$ 各以
        概率 $\frac12$ 取值，证明
        \[
        \E(\eu^{t\varepsilon}) = \cosh t \leq \eu^{t^2/2}
        \qquad (t \in \R)
        \]
        通过逐项比较两个级数（$(2k)! \geq
        2^kk!$）。
  \item 对独立的拉德马赫变量
        $\varepsilon_1, \dots, \varepsilon_n$ 及任意 $s >
        0$，推出：
        \[
        \P\Bigl(\sum_{i=1}^n\varepsilon_i \geq s\Bigr)
        \leq \eu^{-s^2/(2n)} .
        \]
  \item 改写为公平硬币的情形（令 $X_i =
        \frac{1+\varepsilon_i}2$），则 $\P\bigl(\widehat p_n -
        \tfrac12 \geq \delta\bigr) \leq \eu^{-2n\delta^2}$，
        以及带因子 $2$ 的双侧版本。
  \item （霍夫丁引理）设 $X \in \intcc01$ 且 $\E X =
        p$，记 $\psi(t) = \ln\E(\eu^{tX})$，并说明
        $\psi$ 二阶可导，且
        \[
        \psi''(t) = \E_t(X^2) - \E_t(X)^2, \qquad
        \E_t(Y) := \frac{\E(Y\eu^{tX})}{\E(\eu^{tX})},
        \]
        即一个重加权变量的 \emph{方差}，该变量仍取值于
        $\intcc01$ 中，以 $\frac14$ 为其上界
        （\cref{exo:b2:randomvar:6}中的极小性论证），并由泰勒公式
        得出：
        \[
        \E\bigl(\eu^{t(X - p)}\bigr) \leq \eu^{t^2/8} .
        \]
  \item （霍夫丁不等式）对独立的 $X_i \in
        \intcc01$ 且有共同均值 $p$，推出
        \[
        \P\bigl(\abs{\widehat p_n - p} \geq \delta\bigr)
        \leq 2\,\eu^{-2n\delta^2}
        \qquad (\delta > 0).
        \]
  \item 比较切比雪夫速率
        $\frac{p(1-p)}{n\delta^2}$ 与霍夫丁界
        $2\eu^{-2n\delta^2}$：二者各需要何种假设，并
        （大致）从哪个 $n$ 起，
        在 $\delta = 0.03$、$p =
        \frac12$ 时指数界更优？
\end{enumerate}

\medskip
\noindent\textbf{第三部分 --- 民意调查须询问多少人
？} 某次民调独立且均匀地选取 $n$ 位选民；
每人如实作答；$p$ 为真实比例，$\widehat p_n$
为该次民调的结果。
\begin{enumerate}[resume]
  \item 证明：民调结果精确到 $\pm\delta$，
        置信度为 $1 - \alpha$（即\
        $\P(\abs{\widehat p_n - p} \geq \delta) \leq \alpha$），
        只要
        \[
        n \;\geq\; \frac{\ln(2/\alpha)}{2\,\delta^2} .
        \]
  \item 计算所需的 $n$，取标准规格“三个
        百分点、百分之九十五”（$\delta
        = 0.03$，$\alpha = 0.05$）：$n \geq 2050$；若只取一个点，则
        $n \geq 18\,445$。试指出 --- 并解释 --- 一个引人注目的事实：
        答案
        不依赖于总体规模。
  \item 用切比雪夫不等式重做第 13 题（$V(X_1) = p(1-p)
        \leq \frac14$）：$n \geq \frac1{4\alpha\delta^2} =
        5556$（取三个百分点时）。注意，
        \emph{不放回}抽样只会更有利
        （\cref{exo:b2:randomvar:5}：方差乘以
        $\frac{N-n}{N-1}$）。
  \item （预测选举）某候选人的真实比例为 $p
        = 0.52$。至少需要询问多少选民，才能使
        $\P(\widehat p_n \leq \tfrac12) \leq 0.01$？证明 $n
        \geq \frac{\ln 100}{2\cdot(0.02)^2} \approx 5757$ ---
        预测一场胶着的选举，其代价
        远高于估计一个比例。
  \item 数学\emph{未}涵盖的内容：列出
        所用的建模假设（独立、均匀
        抽样、如实作答、$p$ 固定），并用一小段文字说明
        真实民调误差为何主要来自
        \emph{偏差}（抽样不均匀、无应答等），
        而这类误差不会随 $n$ 增大而减小。
\end{enumerate}

\medskip
\noindent\textbf{第四部分 --- 更精确、更经济。}
\begin{enumerate}[resume]
  \item （均值的中位数：由二阶矩得到指数衰减）
        将 $km$ 个样本的预算分成 $k$ 个相互独立的
        组，每组 $m$ 个；记 $\widehat p^{(1)}, \dots, \widehat
        p^{(k)}$ 为各组的均值，$M$ 为它们的中位数。
        取 $m$ 使每组满足
        $\P(\abs{\widehat p^{(i)} - p} \geq \delta) \leq
        \frac18$（由切比雪夫不等式，$m \geq \frac2{\delta^2}$
        即可）。证明：若 $\abs{M - p} \geq \delta$
        则至少有 $k/2$ 组出错，由此推出
        \[
        \P(\abs{M - p} \geq \delta)
        \leq \binom{k}{\lceil k/2\rceil}\Bigl(\frac18
        \Bigr)^{k/2}
        \leq 2^k\cdot 8^{-k/2} = 2^{-k/2} :
        \]
        仅凭方差这一工具，即可得到
        指数级的集中性。
  \item （佩利–齐格蒙德）对具有二阶矩的 $X \geq 0$，
        证明 $\P(X > 0) \geq \dfrac{\E(X)^2}{\E(X^2)}$
        \emph{（对 $X\mathbf 1_{X>0}$ 使用柯西–施瓦茨不等式）}：这是
        一个反方向的工具 --- 矩同样可以迫使
        事件发生。
  \item （平斯克简版）证明 $I(a) \geq 2\bigl(a -
        \tfrac12\bigr)^2$ 在 $\intoo{\frac12}1$ 上成立 \emph{（此差
        在 $\frac12$ 处二阶无穷小，且
        其二阶导数为 $\frac1{a(1-a)} - 4 \geq
        0$）}：切尔诺夫的精确指数总优于
        霍夫丁给出的二次型指数。
  \item 展开 $I\bigl(\tfrac12 + \delta\bigr) = 2\delta^2 +
        O(\delta^4)$，并结合第 4 题：对于小偏差，
        霍夫丁指数 $2n\delta^2$ 渐近
        \emph{精确} --- 任何方法都无法在多项式因子之外
        超越它。
  \item 列出工具表：对于马尔可夫、切比雪夫、
        \cref{exo:b2:randomvar:9}中的四阶矩界、
        霍夫丁，以及以 $I$ 为指数的切尔诺夫界，各用一行写出
        所需假设、所得衰减速率，
        以及本题中该界最为精确的
        那一问。
\end{enumerate}

\medskip
\noindent\textbf{第五部分 --- 红利。}
\begin{enumerate}[resume]
  \item （检验硬币）一枚硬币或是公平的，或是有偏的，其
        正面概率 $p = 0.55$。抛掷 $n$ 次，并在
        $\widehat p_n > 0.525$ 时宣布“有偏”。证明
        两类错误概率均至多为
        $\eu^{-2n(0.025)^2}$，并证明 $n \geq 3685$ 次
        抛掷即可保证两者均小于 $1\%$。
  \item （稀有事件需要考虑方差的界）设 $p =
        0.01$，取相对精度要求 $\delta =
        p/2 = 0.005$，$\alpha = 0.05$。比较
        霍夫丁界（$n \approx 74\,000$）所要求的样本量，以及
        由真实方差 $p(1-p)$ 得到的切比雪夫界所要求的样本量（$n
        \approx 7920$）。忽略方差的指数界输给朴素的二阶矩
        界。并指出其寓意，
        并说明所缺工具（一个考虑方差的指数界；
        泊松近似，见
        \cref{ch:b2:genfun}）从何而来。
  \item （硬币的强大数律）由 $\sum_n
        2\eu^{-2n\delta^2} < \infty$ 及博雷尔–坎泰利
        （\cref{thm:b2:proba:borelcantelli}），证明
        对独立的硬币抛掷，$\widehat p_n \to p$ 几乎必然成立，
        将几乎必然事件写成
        $\bigcap_j\bigcup_N\bigcap_{n\geq N}
        \{\abs{\widehat p_n - p} < \tfrac1j\}$，如
        \cref{exo:b2:randomvar:9}所示，并由此得出结论。（有界性
        代替了其中所用的四阶矩。）
  \item 综合。用五句话说明梯子上每一级
        （一、二、四阶矩；有界指数界；
        精确指数）各自付出什么、得到什么；为何只需调查 $2050$
        人即足以代表任意规模的国家；并指出其中
        哪一个界将在本科三年级分册中精化为
        中心极限定理的精确常数。
\end{enumerate}
\end{problem}
